% unidades/35-repaso-general.tex
\chapter{🎯 \ruby{総復習}{そうふくしゅう} — Repaso General N4→N3}
\unidadheader{35}{総復習}{全35ユニット}{Consolidar gramática N4→N3 y preparar el JLPT N3}

\begin{tcolorbox}[vocabbox, title={\emoji{pushpin} Puntos Clave}]
\begin{tabularx}{\linewidth}{llX}
\toprule
\textbf{Unidades} & \textbf{Tema} & \textbf{Forma clave} \\ \midrule
01--05 & Repaso N4 / て-form avanzado & てしまう・ておく・てある \\ \midrule
06--09 & Pasiva y causativa & 〜られる・〜させる・〜させられる \\ \midrule
10--15 & Propósito, evidencialidad & ために・ように・はずだ・わけだ \\ \midrule
16--20 & Límites, cambio, relativas & だけ・しか・すぎる・関係節 \\ \midrule
21--24 & Condicionales & たら・ば・なら・と \\ \midrule
25--30 & Opinión, keigo & と思う・といわれている・尊敬語・謙譲語 \\ \midrule
31--35 & Registro y lectura & である体・縮約形・読解 \\
\bottomrule
\end{tabularx}
\end{tcolorbox}

\subsection*{\emoji{gear} Mapa de Condicionales — la comparación clave}

\begin{tcolorbox}[comparacionbox, title={⚖ 4 Condicionales — cuándo usar cuál}]
\begin{tabularx}{\linewidth}{lXX}
\toprule
\textbf{Forma} & \textbf{Uso principal} & \textbf{Ejemplo} \\ \midrule
〜たら & Secuencia, hipótesis, descubrimiento & \ruby{着}{つ}いたら\ruby{電話}{でんわ}する \\
〜ば & Condición general, proverbio, arrepentimiento & \ruby{急}{いそ}がば\ruby{回}{まわ}れ \\
〜なら & Respuesta a situación dada, consejo & \ruby{行}{い}くなら\ruby{傘}{かさ}を\ruby{持}{も}って \\
〜と & Resultado automático / inevitable & ボタンを\ruby{押}{お}すと\ruby{点滅}{てんめつ}する \\
\bottomrule
\end{tabularx}
\end{tcolorbox}

\subsection*{\emoji{gear} Pasiva y Causativa — resumen}
\patron{Voz pasiva — el sujeto recibe la acción}
  {G1: u→areru \quad G2: ru→rareru \quad Pasiva de perjuicio: recibe daño indirecto}
  {\ej{\ruby{先生}{せんせい}に\ruby{褒}{ほ}められた。}{Me elogió el profesor. (pasiva directa)}
   \ej{\ruby{雨}{あめ}に\ruby{降}{ふ}られて\ruby{濡}{ぬ}れた。}{Me mojé porque llovió. (pasiva de perjuicio)}}

\patron{Causativa y causativa-pasiva}
  {G1: u→aseru \quad G2: ru→saseru \quad Causativa-pasiva: 〜させられる}
  {\ej{\ruby{子供}{こども}に\ruby{野菜}{やさい}を\ruby{食}{た}べさせた。}{Hice que el niño comiera verduras. (causativa)}
   \ej{\ruby{残業}{ざんぎょう}させられた。}{Me obligaron a hacer horas extra. (causativa-pasiva)}}

\subsection*{\emoji{gear} Keigo — tabla de referencia}

\begin{tcolorbox}[kanjibox, title={敬語 — Referencia rápida}]
\begin{tabularx}{\linewidth}{lllX}
\toprule
\textbf{Neutro} & \textbf{尊敬語} & \textbf{謙譲語} & \textbf{Significado} \\ \midrule
いる & いらっしゃる & おる & estar \\
いく/くる & いらっしゃる & \ruby{参}{まい}る & ir/venir \\
いう & おっしゃる & \ruby{申}{もう}す & decir \\
する & なさる & いたす & hacer \\
たべる/のむ & めしあがる & いただく & comer/beber \\
みる & ご\ruby{覧}{らん}になる & \ruby{拝見}{はいけん}する & ver \\
しる/おもう & ご\ruby{存知}{ぞんじ} & \ruby{存}{ぞん}じる & saber \\
もらう & — & いただく & recibir \\
\bottomrule
\end{tabularx}
\end{tcolorbox}

\subsection*{\emoji{speech-balloon} Frases de Repaso — una por unidad}
\frase{やってしまった！もうあきらめるしかない。}{¡Metí la pata! Ya no queda más que rendirse.}
\frase{\ruby{彼}{かれ}に\ruby{笑}{わら}われた。それはつらかった。}{Me ridiculizó. Eso fue doloroso.}
\frase{\ruby{健康}{けんこう}のために\ruby{毎日}{まいにち}\ruby{歩}{ある}くようにしている。}{Por salud, intento caminar todos los días.}
\frase{この\ruby{映画}{えいが}はもう\ruby{見}{み}たことがあるはずだ。}{Debería haber visto ya esta película.}
\frase{\ruby{もっと早}{もっとはや}く\ruby{始}{はじ}めればよかった。}{Ojalá hubiera empezado antes.}
\frase{\ruby{先生}{せんせい}はいらっしゃいますか？ご\ruby{覧}{らん}になりましたか？}{¿Está el profesor? ¿Lo vio?}

\begin{tcolorbox}[ejerciciobox, title={\emoji{pencil} Autoevaluación Final}]
\begin{enumerate}
  \item ¿たら, ば, なら o と? «Si vas a Kioto, visita el Kinkakuji.»
  \item Convierte a forma pasiva: 「先生は学生を褒めた。」
  \item Elige: 「明日は雨が降る\_\_\_思います。」(と / が / を)
  \item ¿Cuál es el honorífico de 食べる? ¿Y el humilde?
  \item Escribe en である体: 「この問題は難しいです。解決策が必要です。」
\end{enumerate}
\textbf{Respuestas:}
1) なら (\ruby{京都}{きょうと}に\ruby{行}{い}くなら、\ruby{金閣寺}{きんかくじ}へ) \quad
2) \ruby{学生}{がくせい}は\ruby{先生}{せんせい}に\ruby{褒}{ほ}められた \quad
3) と \quad
4) \ruby{尊敬}{そんけい}: めしあがる / \ruby{謙譲}{けんじょう}: いただく \quad
5) この\ruby{問題}{もんだい}は\ruby{難}{むずか}しい。\ruby{解決策}{かいけつさく}が\ruby{必要}{ひつよう}である。
\end{tcolorbox}

\culturabox{\ruby{継続}{けいぞく}が\ruby{力}{ちから} — El poder de la constancia}{
  Llegar aquí significa haber completado 35 unidades de gramática N4→N3. El japonés es un camino largo, pero cada punto de gramática aprendido abre nuevas puertas. 継続は力なり — la constancia es poder. El JLPT N3 es alcanzable con los conocimientos de este libro; el siguiente paso es abundante práctica de lectura, escucha y conversación. ¡\ruby{頑張}{がんば}ってください！
}

\begin{tcolorbox}[kanjibox, title={\emoji{mahjong} Kanjis del Repaso}]
\begin{tabularx}{\linewidth}{cllXX}
\toprule
\textbf{Kanji} & \textbf{On} & \textbf{Kun} & \textbf{Significado} & \textbf{Vocab} \\ \midrule
\kanjirow{総}{そう}{—}{total/general}{\ruby{総復習}{そうふくしゅう} / \ruby{総合}{そうごう}}
\kanjirow{復}{ふく}{—}{repetir/volver}{\ruby{復習}{ふくしゅう} / \ruby{回復}{かいふく}}
\kanjirow{習}{しゅう}{\ruby{なら}{う}}{aprender/practicar}{\ruby{練習}{れんしゅう} / \ruby{習慣}{しゅうかん}}
\kanjirow{継}{けい}{\ruby{つ}{ぐ}}{continuar}{\ruby{継続}{けいぞく} / \ruby{継承}{けいしょう}}
\bottomrule
\end{tabularx}
\end{tcolorbox}
